% Môn: Vật lý
% Lớp: 10
% Chương: Chương I. Mở đầu
% Bài: L10C1 Bài 2. Các quy tắc an toàn trong phòng thực hành Vật lí · Phân tích mối nguy và đánh giá mức độ rủi ro
% Số câu: 185
% App đọc file này trực tiếp — sửa rồi Commit trên GitHub, bấm Đồng bộ GitHub .tex

% L10C1 Bài 2. Các quy tắc an toàn trong phòng thực hành Vật lí



% ===== Câu 104 =====
% ID: L10C1B2-29-DS
% Mức: NB
\dangbt{Phân tích mối nguy và đánh giá mức độ rủi ro}
\begin{ex}
% ID: L10C1B2-29-DS
% Mức: NB
Trong tình huống khi khu vực làm việc bị ướt đối với đánh giá rủi ro, hãy đánh giá các phát biểu sau.
\choiceTF
{\True Cần xác định mối nguy, ước lượng khả năng xảy ra và mức độ hậu quả trước khi chọn biện pháp kiểm soát.}
{Có thể chỉ xem xét hậu quả mà bỏ qua khả năng xảy ra nếu thao tác diễn ra nhanh.}
{\True Phải báo cho giáo viên hoặc người phụ trách khi phát hiện nguy cơ vượt khả năng kiểm soát.}
{\True Chỉ tiếp tục thí nghiệm sau khi nguyên nhân nguy hiểm đã được loại bỏ hoặc kiểm soát.}
\loigiai{
\begin{itemchoice}
\item (Đúng) Cần xác định mối nguy, ước lượng khả năng xảy ra và mức độ hậu quả trước khi chọn biện pháp kiểm soát. (Vì): vì đây là biện pháp an toàn của dạng đánh giá rủi ro; b sai vì hành vi “chỉ xem xét hậu quả mà bỏ qua khả năng xảy ra” vẫn nguy hiểm; c và d đúng do sự cố phải được báo cáo, cô lập và kiểm soát trước khi tiếp tục
\item (Sai) Có thể chỉ xem xét hậu quả mà bỏ qua khả năng xảy ra nếu thao tác diễn ra nhanh.
\item (Đúng) Phải báo cho giáo viên hoặc người phụ trách khi phát hiện nguy cơ vượt khả năng kiểm soát.
\item (Đúng) Chỉ tiếp tục thí nghiệm sau khi nguyên nhân nguy hiểm đã được loại bỏ hoặc kiểm soát.
\end{itemchoice}
}
\end{ex}

% ===== Câu 106 =====
% ID: L10C1B2-30-DS
% Mức: NB
\begin{ex}
% ID: L10C1B2-30-DS
% Mức: NB
Trong tình huống khi cần rời bàn thí nghiệm đối với đánh giá rủi ro, hãy đánh giá các phát biểu sau.
\choiceTF
{\True Cần xác định mối nguy, ước lượng khả năng xảy ra và mức độ hậu quả trước khi chọn biện pháp kiểm soát.}
{Có thể chỉ xem xét hậu quả mà bỏ qua khả năng xảy ra nếu thao tác diễn ra nhanh.}
{\True Phải báo cho giáo viên hoặc người phụ trách khi phát hiện nguy cơ vượt khả năng kiểm soát.}
{\True Chỉ tiếp tục thí nghiệm sau khi nguyên nhân nguy hiểm đã được loại bỏ hoặc kiểm soát.}
\loigiai{
\begin{itemchoice}
\item (Đúng) Cần xác định mối nguy, ước lượng khả năng xảy ra và mức độ hậu quả trước khi chọn biện pháp kiểm soát. (Vì): vì đây là biện pháp an toàn của dạng đánh giá rủi ro; b sai vì hành vi “chỉ xem xét hậu quả mà bỏ qua khả năng xảy ra” vẫn nguy hiểm; c và d đúng do sự cố phải được báo cáo, cô lập và kiểm soát trước khi tiếp tục
\item (Sai) Có thể chỉ xem xét hậu quả mà bỏ qua khả năng xảy ra nếu thao tác diễn ra nhanh.
\item (Đúng) Phải báo cho giáo viên hoặc người phụ trách khi phát hiện nguy cơ vượt khả năng kiểm soát.
\item (Đúng) Chỉ tiếp tục thí nghiệm sau khi nguyên nhân nguy hiểm đã được loại bỏ hoặc kiểm soát.
\end{itemchoice}
}
\end{ex}

% ===== Câu 108 =====
% ID: L10C1B2-31-DS
% Mức: TH
\begin{ex}
% ID: L10C1B2-31-DS
% Mức: TH
Trong tình huống khi có người không đeo phương tiện bảo hộ đối với đánh giá rủi ro, hãy đánh giá các phát biểu sau.
\choiceTF
{\True Cần xác định mối nguy, ước lượng khả năng xảy ra và mức độ hậu quả trước khi chọn biện pháp kiểm soát.}
{Có thể chỉ xem xét hậu quả mà bỏ qua khả năng xảy ra nếu thao tác diễn ra nhanh.}
{\True Phải báo cho giáo viên hoặc người phụ trách khi phát hiện nguy cơ vượt khả năng kiểm soát.}
{\True Chỉ tiếp tục thí nghiệm sau khi nguyên nhân nguy hiểm đã được loại bỏ hoặc kiểm soát.}
\loigiai{
\begin{itemchoice}
\item (Đúng) Cần xác định mối nguy, ước lượng khả năng xảy ra và mức độ hậu quả trước khi chọn biện pháp kiểm soát. (Vì): vì đây là biện pháp an toàn của dạng đánh giá rủi ro; b sai vì hành vi “chỉ xem xét hậu quả mà bỏ qua khả năng xảy ra” vẫn nguy hiểm; c và d đúng do sự cố phải được báo cáo, cô lập và kiểm soát trước khi tiếp tục
\item (Sai) Có thể chỉ xem xét hậu quả mà bỏ qua khả năng xảy ra nếu thao tác diễn ra nhanh.
\item (Đúng) Phải báo cho giáo viên hoặc người phụ trách khi phát hiện nguy cơ vượt khả năng kiểm soát.
\item (Đúng) Chỉ tiếp tục thí nghiệm sau khi nguyên nhân nguy hiểm đã được loại bỏ hoặc kiểm soát.
\end{itemchoice}
}
\end{ex}

% ===== Câu 109 =====
% ID: L10C1B2-31-TL
% Mức: TH
\begin{bt}
% ID: L10C1B2-31-TL
% Mức: TH
Trình bày nguyên tắc cốt lõi khi làm việc với đánh giá rủi ro và giải thích vì sao phải tuân thủ.
\loigiai{
Nguyên tắc cốt lõi là xác định mối nguy, ước lượng khả năng xảy ra và mức độ hậu quả trước khi chọn biện pháp kiểm soát. Phải tuân thủ vì giúp ưu tiên xử lí nguy cơ có rủi ro cao; đồng thời luôn dừng và báo người phụ trách khi điều kiện không bảo đảm.
}
\end{bt}

% ===== Câu 110 =====
% ID: L10C1B2-31-TLN
% Mức: NB
\begin{ex}
% ID: L10C1B2-31-TLN
% Mức: NB
Một quy trình về đánh giá rủi ro gồm: (1) nhận diện nguy cơ; (2) xác định mối nguy, ước lượng khả năng xảy ra và mức độ hậu quả trước khi chọn biện pháp kiểm soát; (3) chỉ xem xét hậu quả mà bỏ qua khả năng xảy ra; (4) báo giáo viên khi có bất thường; (5) chỉ tiếp tục khi nguy cơ đã được kiểm soát. Có bao nhiêu bước phù hợp quy tắc an toàn? Chỉ ghi số.
\shortans{4}
\loigiai{
Các bước (1), (2), (4), (5) phù hợp; bước (3) không an toàn. Có 4 bước phù hợp.
}
\end{ex}

% ===== Câu 111 =====
% ID: L10C1B2-32-DS
% Mức: TH
\begin{ex}
% ID: L10C1B2-32-DS
% Mức: TH
Trong tình huống khi kết quả đo khác dự kiến đối với đánh giá rủi ro, hãy đánh giá các phát biểu sau.
\choiceTF
{\True Cần xác định mối nguy, ước lượng khả năng xảy ra và mức độ hậu quả trước khi chọn biện pháp kiểm soát.}
{Có thể chỉ xem xét hậu quả mà bỏ qua khả năng xảy ra nếu thao tác diễn ra nhanh.}
{\True Phải báo cho giáo viên hoặc người phụ trách khi phát hiện nguy cơ vượt khả năng kiểm soát.}
{\True Chỉ tiếp tục thí nghiệm sau khi nguyên nhân nguy hiểm đã được loại bỏ hoặc kiểm soát.}
\loigiai{
\begin{itemchoice}
\item (Đúng) Cần xác định mối nguy, ước lượng khả năng xảy ra và mức độ hậu quả trước khi chọn biện pháp kiểm soát. (Vì): vì đây là biện pháp an toàn của dạng đánh giá rủi ro; b sai vì hành vi “chỉ xem xét hậu quả mà bỏ qua khả năng xảy ra” vẫn nguy hiểm; c và d đúng do sự cố phải được báo cáo, cô lập và kiểm soát trước khi tiếp tục
\item (Sai) Có thể chỉ xem xét hậu quả mà bỏ qua khả năng xảy ra nếu thao tác diễn ra nhanh.
\item (Đúng) Phải báo cho giáo viên hoặc người phụ trách khi phát hiện nguy cơ vượt khả năng kiểm soát.
\item (Đúng) Chỉ tiếp tục thí nghiệm sau khi nguyên nhân nguy hiểm đã được loại bỏ hoặc kiểm soát.
\end{itemchoice}
}
\end{ex}

% ===== Câu 112 =====
% ID: L10C1B2-32-TL
% Mức: TH
\begin{bt}
% ID: L10C1B2-32-TL
% Mức: TH
Phân tích hành vi “chỉ xem xét hậu quả mà bỏ qua khả năng xảy ra” và đề xuất cách thay thế an toàn.
\loigiai{
Hành vi “chỉ xem xét hậu quả mà bỏ qua khả năng xảy ra” có thể gây tai nạn vì làm mất kiểm soát nguồn nguy hiểm. Cần dừng thao tác, cô lập nguy cơ và thay bằng biện pháp: xác định mối nguy, ước lượng khả năng xảy ra và mức độ hậu quả trước khi chọn biện pháp kiểm soát.
}
\end{bt}

% ===== Câu 117 =====
% ID: L10C1B2-35-TL
% Mức: VD
\begin{bt}
% ID: L10C1B2-35-TL
% Mức: VD
So sánh biện pháp phòng ngừa và biện pháp ứng phó đối với nguy cơ thuộc nhóm đánh giá rủi ro.
\loigiai{
Phòng ngừa diễn ra trước sự cố, tập trung vào kiểm tra, loại bỏ hoặc giảm nguy cơ. Ứng phó diễn ra khi sự cố đã hoặc sắp xảy ra, gồm dừng hoạt động, cô lập nguồn nguy hiểm, báo cáo và sơ cứu. Hai nhóm biện pháp bổ sung cho nhau.
}
\end{bt}

% ===== Câu 119 =====
% ID: L10C1B2-36-TL
% Mức: VDC
\begin{bt}
% ID: L10C1B2-36-TL
% Mức: VDC
Đề xuất cách tổ chức nhóm học sinh để thực hiện nhiệm vụ về đánh giá rủi ro an toàn và có thể kiểm tra được.
\loigiai{
Phân công trưởng nhóm kiểm tra điều kiện, người thao tác, người quan sát an toàn và người ghi chép. Dùng bảng kiểm, xác nhận trước mỗi bước, quy định tín hiệu dừng và báo ngay bất thường.
}
\end{bt}

% ===== Câu 161 =====
% ID: L10C1B2-67-TN
% Mức: NB
\begin{ex}
% ID: L10C1B2-67-TN
% Mức: NB
Trong lúc thay đổi cách bố trí dụng cụ liên quan đến đánh giá rủi ro, hành động nào an toàn nhất?
\choice
{tiếp tục thao tác rồi báo cáo sau}
{nhờ bạn khác làm thay mà không kiểm tra điều kiện an toàn}
{\True xác định mối nguy, ước lượng khả năng xảy ra và mức độ hậu quả trước khi chọn biện pháp kiểm soát}
{chỉ xem xét hậu quả mà bỏ qua khả năng xảy ra}
\loigiai{
Để đánh giá rủi ro một cách an toàn, trước hết cần xác định các mối nguy tiềm ẩn, sau đó ước lượng khả năng xảy ra của các mối nguy đó và mức độ hậu quả có thể xảy ra. Dựa trên những thông tin này, ta mới có thể lựa chọn các biện pháp kiểm soát phù hợp và hiệu quả nhất, ưu tiên xử lý những rủi ro có mức độ cao
}
\end{ex}

% ===== Câu 162 =====
% ID: L10C1B2-68-TN
% Mức: TH
\begin{ex}
% ID: L10C1B2-68-TN
% Mức: TH
Sau khi hoàn thành phép đo liên quan đến đánh giá rủi ro, hành động nào an toàn nhất?
\choice
{nhờ bạn khác làm thay mà không kiểm tra điều kiện an toàn}
{\True xác định mối nguy, ước lượng khả năng xảy ra và mức độ hậu quả trước khi chọn biện pháp kiểm soát}
{chỉ xem xét hậu quả mà bỏ qua khả năng xảy ra}
{tiếp tục thao tác rồi báo cáo sau}
\loigiai{
Hành động đúng là “xác định mối nguy, ước lượng khả năng xảy ra và mức độ hậu quả trước khi chọn biện pháp kiểm soát” vì giúp ưu tiên xử lí nguy cơ có rủi ro cao. Chọn B.
}
\end{ex}

% ===== Câu 163 =====
% ID: L10C1B2-70-TN
% Mức: TH
\begin{ex}
% ID: L10C1B2-70-TN
% Mức: TH
Khi khu vực làm việc bị ướt liên quan đến đánh giá rủi ro, hành động nào an toàn nhất?
\choice
{chỉ xem xét hậu quả mà bỏ qua khả năng xảy ra}
{tiếp tục thao tác rồi báo cáo sau}
{nhờ bạn khác làm thay mà không kiểm tra điều kiện an toàn}
{\True xác định mối nguy, ước lượng khả năng xảy ra và mức độ hậu quả trước khi chọn biện pháp kiểm soát}
\loigiai{
Hành động đúng là “xác định mối nguy, ước lượng khả năng xảy ra và mức độ hậu quả trước khi chọn biện pháp kiểm soát” vì giúp ưu tiên xử lí nguy cơ có rủi ro cao. Chọn D.
}
\end{ex}

% ===== Câu 164 =====
% ID: L10C1B2-71-TN
% Mức: TH
\begin{ex}
% ID: L10C1B2-71-TN
% Mức: TH
Khi cần rời bàn thí nghiệm liên quan đến đánh giá rủi ro, hành động nào an toàn nhất?
\choice
{tiếp tục thao tác rồi báo cáo sau}
{nhờ bạn khác làm thay mà không kiểm tra điều kiện an toàn}
{\True xác định mối nguy, ước lượng khả năng xảy ra và mức độ hậu quả trước khi chọn biện pháp kiểm soát}
{chỉ xem xét hậu quả mà bỏ qua khả năng xảy ra}
\loigiai{
Hành động đúng là “xác định mối nguy, ước lượng khả năng xảy ra và mức độ hậu quả trước khi chọn biện pháp kiểm soát” vì giúp ưu tiên xử lí nguy cơ có rủi ro cao. Chọn C.
}
\end{ex}

% ===== Câu 165 =====
% ID: L10C1B2-72-TN
% Mức: VD
\begin{ex}
% ID: L10C1B2-72-TN
% Mức: VD
Khi có người không đeo phương tiện bảo hộ liên quan đến đánh giá rủi ro, hành động nào an toàn nhất?
\choice
{nhờ bạn khác làm thay mà không kiểm tra điều kiện an toàn}
{\True xác định mối nguy, ước lượng khả năng xảy ra và mức độ hậu quả trước khi chọn biện pháp kiểm soát}
{chỉ xem xét hậu quả mà bỏ qua khả năng xảy ra}
{tiếp tục thao tác rồi báo cáo sau}
\loigiai{
Hành động đúng là “xác định mối nguy, ước lượng khả năng xảy ra và mức độ hậu quả trước khi chọn biện pháp kiểm soát” vì giúp ưu tiên xử lí nguy cơ có rủi ro cao. Chọn B.
}
\end{ex}

% ===== Câu 166 =====
% ID: L10C1B2-74-TN
% Mức: VD
\begin{ex}
% ID: L10C1B2-74-TN
% Mức: VD
Khi giáo viên yêu cầu dừng thao tác liên quan đến đánh giá rủi ro, hành động nào an toàn nhất?
\choice
{chỉ xem xét hậu quả mà bỏ qua khả năng xảy ra}
{tiếp tục thao tác rồi báo cáo sau}
{nhờ bạn khác làm thay mà không kiểm tra điều kiện an toàn}
{\True xác định mối nguy, ước lượng khả năng xảy ra và mức độ hậu quả trước khi chọn biện pháp kiểm soát}
\loigiai{
Hành động đúng là “xác định mối nguy, ước lượng khả năng xảy ra và mức độ hậu quả trước khi chọn biện pháp kiểm soát” vì giúp ưu tiên xử lí nguy cơ có rủi ro cao. Chọn D.
}
\end{ex}

